\section{The Nucleus Spectrum Index}
\label{sec:index}

In this section, we describe how \specnd{} stores its output as a queryable index, \nsi, whose size scales with the graph's clique structure rather than its clique count.
%
We first describe the layout, then the query procedure, and finally the size.

\stitle{What the index stores.}
A spectrum cannot be stored as values, since one row can hold more than $10^{12}$ $r$-cliques.
%
\nsi{} stores structure instead.
%
It records the class of every vertex, the vertex sets of the mergeable regions, and, for each pattern, its class multiset, its clique bound $\cliqb(P)$, and its boundary-cell core value.
%
For each cell above the boundary it stores only a \textit{residue dictionary}, the core values of the patterns that were not certified at that cell.
%
Certified values are not stored at all; they are reconstructed by arithmetic at query time.

\stitle{Answering a query.}
Given an $r$-clique $R$ and a cell $s$, \nsi{} returns $\corev{r,s}(R)$ as follows.
%
It maps the vertices of $R$ to their classes and forms the pattern; if the pattern is present, it walks the chain from the boundary cell, and at each step a certified pattern yields the closed form $\binom{\cliqb(P)-r}{s-r}$ by Theorem~\ref{thm:chain} while an uncertified pattern is read from that cell's residue dictionary.
%
Because certification is absorbing, the walk stops at the first certified step and returns the closed form directly for all higher cells.
%
If the pattern is absent, $R$ lies in no active region; if some mergeable region contains all of $R$, the answer is the closed form $\binom{|M|-r}{s-r}$ by Theorem~\ref{thm:noearly}, and otherwise $R$ is in no $s$-clique and the answer is $0$.
%
The pattern case and the mergeable case are disjoint by Theorem~\ref{thm:quotient}, so the lookup is unambiguous.

\begin{example}
\label{exa:query}
To answer $\corev{3,6}(\{v_3,v_4,v_5\})$ on Figure~\ref{fig:running}, \nsi{} forms the pattern of $\{v_3,v_4,v_5\}$, reads its boundary value $\corev{3,4}=2$, and finds it certified at $s{=}5$; certification is absorbing, so it returns $\binom{2}{6-3}=\binom{2}{3}=0$ immediately without consulting the residue dictionaries.
\end{example}

\stitle{Size.}
The index holds one record per pattern, one per mergeable region, and one residue entry per uncertified pattern per cell.
%
When most patterns certify, the residue dictionaries are nearly empty and the index size is dominated by the pattern records, which are independent of the number of $r$-cliques.
%
On \daWIT, the whole four-cell spectrum of $338$ million $r$-cliques occupies $14.5$MB, that is, $0.045$ bytes per $r$-clique; a per-cell listing of the same spectrum would store each $r$-clique once per cell and require several orders of magnitude more space.
%
On graphs with a large residue the dictionaries dominate, and the index is larger but still holds the whole spectrum in tens of megabytes (Section~\ref{sec:exp}).
